\chapter[{[GEF]} Gefahrenzonen, Gefahrentransporte,
Gefahrengutunfall]{[GEF]\\Gefahrenzonen, \\Gefahrentransporte,
\\Gefahrengutunfall}
\label{topic:gefahrengutunfall}

\minitoc

% TODO : Layout-Anpassung
\TODO{GABS}{yyyy-mm-dd}{Layout-Anpassung notwendig.}{Dieser Teil entspricht nicht dem Standard-Layout!}


\noindent\Parallel{
% \includegraphics[width=\linewidth]{./images/noauth/AASdev20090228-img149.jpg}
\includegraphics[width=\linewidth]{./icons/under-construction.png}
}{
% \includegraphics[width=\linewidth]{./images/noauth/AASdev20090228-img150.jpg}
\includegraphics[width=\linewidth]{./icons/under-construction.png}
}

\TODO{GABS}{2010-04-27}{Bilder 2x Gefahrengutunfall/LKW}{}

\section{Gefahrenzonen}

\Layout{0}{Mögliche Gefahren für den Sanitäter}{
\begin{itemize}
    \item Brand bzw. Explosionsgefahr
    \item Entweichen unter Druck stehender Gase
    \item Giftigkeit, Ansteckungsgefahr
    \item \"Atzwirkung
    \item Radioaktivität
    \item Umweltgefährdung
\end{itemize}
}{
\begin{itemize}
    \item Brand bzw. Explosionsgefahr
    \item Entweichen unter Druck stehender Gase
    \item Giftigkeit, Ansteckungsgefahr
    \item \"Atzwirkung
    \item Radioaktivität
    \item Umweltgefährdung
\end{itemize}
}{}{}{}




\Layout{0}{Gefahrenquellen können entstehen bei}{
\begin{itemize}
    \item Unfälle bei der Herstellung
    \item Transportunfälle
    \item Unfälle bei der Verwendung von gefährlichen Stoffen
\end{itemize}
}{
\begin{itemize}
    \item Unfälle bei der Herstellung
    \item Transportunfälle
    \item Unfälle bei der Verwendung von gefährlichen Stoffen   
\end{itemize}
}{}{}{}



\noindent\begin{tabulary}{\linewidth}{LL}\toprule
\multicolumn{2}{c}{\TabHeading{Andere Gefahrenquellen:}}
\\\midrule
Landwirtschaft
 &
\ce{CO2} in Gärkellern, Siloanlagen, Brunenschächte

Pflanzenschutzmittel
\\\addlinespace
Klärgruben
 &
\ce{CO2}, Methan, Schwefelwasserstoff
\\\addlinespace
Obstlagerhallen
 &
wenig \ce{O2}, viel \ce{CO2} oder \ce{N2} halten Obst frisch
\\\addlinespace
Brände in Wirtschaftsgebäuden
 &
extrem gefährliche Rauchgase
\\\addlinespace
Tiefgaragen und Tunnels
 &
erhöhte \ce{CO}- und \ce{CO2}-Werte
\\\addlinespace
Flüssiggas
 &
Campingplätze, feste und mobile Küchen
\\\addlinespace
Pools, Schwimmbäder
 &
Chlor (Wasserdesinfektion)
\\\addlinespace
Industrie
 &
Kühlmittel in großen Kälteanlagen (Ammoniak)
\\\addlinespace
Speziallabors, Einrichtungen mit und zur Verarbeitung von infektiösen
Abfällen oder Tierkadaver
 &
Infektionsgefahr
\\\addlinespace
Nuklearmedizinische Institute, Industrie, Forschungseinrichtungen
 &
Radioaktivität
\\\bottomrule
\end{tabulary}






















\section{Kennzeichnung von gefährlichen Stoffen und
Gefahrenbereichen}
% TODO : Layout-Anpassung
\TODO{GABS}{yyyy-mm-dd}{Layout-Anpassung notwendig.}{Dieser Teil entspricht nicht dem Standard-Layout!}

\begin{center}
\begin{tabulary}{\textwidth}{LLLL}
\toprule
\multicolumn{2}{c}{\TabHeading{Hinweis auf Gefahr}} 
&
\multicolumn{2}{c}{\TabHeading{Warnzeichen}}
\\\midrule
% Unhandled or unsupported graphics:
\includegraphics[width=1.475cm,height=1.452cm]{./images/noauth/gefahr-explosion.png}
 &
Explosionsgefährlich
 &
\includegraphics[width=1.541cm,height=1.426cm]{./images/noauth/warnung-allgemein.png}
 &
Allgemeine Gefahr
\\
\includegraphics[width=1.452cm,height=1.414cm]{./images/noauth/gefahr-giftig.png}
 &
Giftig
 &
\includegraphics[width=1.673cm,height=1.414cm]{./images/noauth/warnung-radioaktiv.png}
 &
Radioaktive Stoffe, ionisierende Strahlen
\\
\includegraphics[width=1.497cm,height=1.611cm]{./images/noauth/gefahr-hochentzuendlich.png}
 &
Hochentzündlich
 &
\includegraphics[width=1.696cm,height=1.412cm]{./images/noauth/warnung-gasflaschen.png}
 &
Gasflaschen
\\
\includegraphics[width=1.475cm,height=1.435cm]{./images/noauth/gefahr-reizend.png}
 &
Reizend
 &
\includegraphics[width=1.673cm,height=1.414cm]{./images/noauth/warnung-laser.png}
 &
Laserstrahlen
\\
\includegraphics[width=1.488cm,height=1.599cm]{./images/noauth/gefahr-aetzend.png}
 &
\"Atzend
 &
 &
\\\hline
\end{tabulary}
\end{center}

\subsection{Transport gefährlicher Güter}

Kennzeichnungspflicht besteht erst dann, wenn die transportierte Menge
eines Stoffes die gesetzlich erlaubte Mindestmenge überschreitet.
Diese Mindestmenge gilt für jeden einzelnen Stoff und nicht für die
Gesamtmenge.

Es können also erhebliche Mengen gefährlicher Güter ohne
Kennzeichnung des Fahrzeuges transportiert werden. 

\subsection{Gefahrzettel}

Bei Stückguttransporten wird jedes Versandstück gekennzeichnet. An
Fahrzeugen, die kein Stückgut transportieren, müssen Gefahrenzettel
an den Längsseiten und an der Rückseite angebracht werden.

\begin{flushleft}
\begin{tabular}{m{2.cm}m{3.5cm}m{2.cm}m{3.5cm}}
 Beispiele
 &
 &
 &
\\
\includegraphics[width=1.659cm,height=1.885cm]{./images/noauth/gefahrzettel-feuergefaehrlich-fluessig.png}
 &
Feuergefährlich \par(flüssige Stoffe)
 &
\includegraphics[width=1.856cm,height=2.096cm]{./images/noauth/gefahrzettel-selbstentzuendlich.png}
 &
Selbstentzündlich
\\
\includegraphics[width=1.835cm,height=1.902cm]{./images/noauth/gefahrzettel-feuergefaehrlich-fest.png}
 &
Feuergefährlich \par(feste Stoffe)
 &
\includegraphics[width=1.909cm,height=2.162cm]{./images/noauth/gefahrzettel-enzuendlich-wasser.png}
 &
Entzündbare Gase bei Berührung mit \Es{Wasser}
\\
\end{tabular}
\end{flushleft}

\subsection{Gefahrentafel}

\begin{table}[H]
\Captionof{table}{Kennzeichnung durch Gefahrentafel.}

\begin{tabular}{m{3.225cm}m{3.225cm}m{3.225cm}m{3.225cm}}
\multicolumn{2}{m{6.65cm}}{\centering Allgemeine Kennzeichnung
(Sammeltransporte)} &
\multicolumn{2}{m{6.65cm}}{\centering Spezielle
Kennzeichnung} 
\\\midrule
\includegraphics[width=\linewidth]{./images/khd/warntafel-allgemein.png}
 &
 &
\includegraphics[width=\linewidth]{./images/khd/warntafel-un-60-1710.png}
 &
Gefahrnummer\par (\T{Kemler-Nummer})

Stoffnummer\par (\T{UN-Nummer})
\\\bottomrule
\end{tabular}
\end{table}


\subsection{Gefahrnummer}

\begin{table}[H]
\Captionof{table}{Gefahrennummer.}

\begin{tabular}{cp{11cm}}
\addlinespace
 2 & Entweichen von Gas durch Druck oder chemische Reaktion
\\\addlinespace
 3 & Entzündbarkeit von Flüssigkeiten (Dämpfen) und Gasen
\\\addlinespace
 4 & Entzündbarkeit fester Stoffe
\\\addlinespace
 5 & Oxidierende (brandfördernde) Wirkung
\\\addlinespace
 6 & Giftig
\\\addlinespace
 7 & Radioaktivität
\\\addlinespace
 8 & \"Atzwirkung
\\\addlinespace
 9 & Gefahr einer spontan heftigen Reaktion
\\\addlinespace\bottomrule
\end{tabular}
\end{table}


Die Verdopplung einer Ziffer weist auf die Zunahme der entsprechenden 
Gefahr hin.

X  vor der Nummer zur Kennzeichnung der Gefahr ${\rightarrow}$ Stoff
reagiert in  gefährlicher Weise mit Wasser

Wenn die Gefahr eines Stoffes ausreichend von einer einzigen Ziffer 
angegeben werden kann, wird dieser Ziffer eine Null angefügt.

Folgende Ziffernkombinationen haben eine besondere Bedeutung

\begin{table}[H]
\Captionof{table}{Spezielle Ziffernkombinationen bei der Gefahrennummer.}

\begin{tabular}{cp{11cm}}
\addlinespace
22 & tiefgekühltes Gas
\\\addlinespace
X323 & entzündbarer flüssiger Stoff, der mit Wasser gefährlich reagiert, wobei entzündbare Gase entweichen
\\\addlinespace
X333 & selbstentzündliche Flüssigkeit, die mit Wasser gefährlich reagiert
\\\addlinespace
X423 & entzündbarer fester Stoff, der mit Wasser gefährlich reagiert, wobei brennbare Gase entweichen
\\\addlinespace
44 & entzündbarer fester Stoff, der sich bei erhöhter Temperatur in Geschmolzenem Zustand befindet
\\\addlinespace
539 & entzündbares organisches Peroxid
\\\addlinespace
90  & verschiedene gefährlich Stoffe
\\\addlinespace\bottomrule
\end{tabular}
\end{table}

















\section{Allgemeine Einsatzrichtlinie}
% TODO : Layout-Anpassung
\TODO{GABS}{yyyy-mm-dd}{Layout-Anpassung notwendig.}{Dieser Teil entspricht nicht dem Standard-Layout!}



\begin{description}
    \item[G]efahr erkennen
    \item[A]bsperrmaßnahmen durchführen
    \item[S]pezialkräfte anfordern
\end{description}







%\begin{gWideParEnv}
%\begin{tabular}{p{0.2\linewidth}p{0.5\linewidth}p{0.3\linewidth}}
\noindent\begin{tabular}{p{3cm}p{\linewidth - 7.6cm}p{3.1cm}}
\toprule
\noindent\TabHeading{{\color{ubuntured}G}efahr erkennen}
 &
\noindent Die Suche nach Warntafeln ist bei umgestürzten Fahrzeugen erschwert
 &
~

\includegraphics[width=3cm]{./images/noauth/gefahrgut-unfall.jpg}
\\
\noindent\TabHeading{{\color{ubuntured}A}bsperr\-maßnahmen
durch\-führen} &
Die Absperrung sollte großräumig  sein, genügender Platz für
die technischen Hilfeleistungen muss einkalkuliert werden. Beachte sich
ändernde Situationen  (Windrichtung, Niederschlag)
 &
~

\includegraphics[width=3cm]{./images/noauth/gefahrgut-absperrung.png}
\\
\noindent\TabHeading{{\color{ubuntured}S}pezial\-kräfte anfordern}
 &
Informationen angeben (Kemmler-Nummer, etc.)
 &
~

\includegraphics[width=3cm]{./images/noauth/gefahrgut-feuerwehr.jpg}
\\\bottomrule
\end{tabular}
%\end{gWideParEnv}


\subsection{Gefahr erkennen}


\subsection{Absperrmaßnahmen durchführen}
% TODO : Layout-Anpassung
\TODO{GABS}{yyyy-mm-dd}{Layout-Anpassung notwendig.}{Dieser Teil entspricht nicht dem Standard-Layout!}

\TakeNotesLarge

\begin{table}[h]
\begin{gWideParEnv}
\Captionof{table}{Bilderserie \SlideHeading{Absperrung}.}

\begin{tabularx}{\linewidth}[H]{XXX}
\includegraphics[width=\linewidth]{./images/khd/gef-wind-rauch1.jpg}
\Captionof{figure}{\AuthNotesPicLegalOk{}{}Ein Auto brennt und alle schauen
fasziniert zu \ldots \SourcePicture{Gabriel}} 
&
\includegraphics[width=\linewidth]{./images/khd/gef-wind-rauch2.jpg}
\Captionof{figure}{\AuthNotesPicLegalOk{}{}\ldots doch dann kommt Wind auf \ldots  
\SourcePicture{Gabriel}} 
&
\includegraphics[width=\linewidth]{./images/khd/gef-wind-rauch.jpg}
\Captionof{figure}{\AuthNotesPicLegalOk{}{}Und oh Schreck! Die Absperrung
verliert ihren Zweck!  \SourcePicture{Gabriel}} 
\\\bottomrule
\end{tabularx}

\Fazit{Die Moral aus der Geschicht': Vertrau' dem lauen Klima nicht!}
\end{gWideParEnv}
\end{table}

\TakeNotesLarge


\subsection{Spezialkräfte anfordern}
% TODO : Layout-Anpassung
\TODO{GABS}{yyyy-mm-dd}{Layout-Anpassung notwendig.}{Dieser Teil entspricht nicht dem Standard-Layout!}


\TakeNotesLarge
% TODO : Neufassung
\TODO{GABS}{yyyy-mm-dd}{Fehlender Text.}{}

\nocite{Cicero911CommissionReport}